\section{Introducción}

\subsection{Objetivos del trabajo práctico}

El presente trabajo práctico tuvo como objetivo familiarizarse con los instrumentos utilizados para medir magnitudes en circuitos de corriente alterna. También se aplicaron los conocimientos teóricos de corriente alterna a un circuito real, estudiando la relación entre la potencia activa, reactiva y aparente. Finalmente, se analizó el método de corrección del factor de potencia mediante bancos de capacitores y el efecto producido por la incorporación de un núcleo en un inductor conectado al circuito.

\subsection{Circuito analizado}

El esquemático correspondiente al circuito estudiado se muestra en la Figura~\ref{fig:circuito}. Este consistió en una fuente de tensión alterna variable, representada como $V_S$; un amperímetro, indicado con la letra A, utilizado para medir la corriente eficaz total consumida por el circuito; un vatímetro, indicado con la letra W, utilizado para medir la potencia activa entregada a la carga; y un voltímetro, indicado con la letra V, utilizado para medir la tensión eficaz aplicada sobre la carga.

La carga estuvo formada por una rama resistiva compuesta por dos lámparas incandescentes, modelada mediante una resistencia equivalente $R$; una rama inductiva compuesta por una bobina, modelada mediante una inductancia $L$ y su resistencia en serie asociada $R_L$; y un banco de capacitores conectable en paralelo, representado mediante una capacitancia $C$. La conexión de cada rama se controló mediante las llaves $S_R$, $S_L$ y $S_C$, donde $S_R$ permitió conectar o desconectar la rama resistiva, $S_L$ la rama inductiva y $S_C$ la rama capacitiva.


Durante la experiencia se trabajó con dos bobinas distintas, denominadas $L_1$ y $L_2$. La bobina $L_1$ fue la utilizada originalmente para las primeras mediciones. Sin embargo, en los ensayos con núcleos metálicos se observó que algunas respuestas experimentales diferían considerablemente de lo esperado a partir de las estimaciones teóricas. Por este motivo, se decidió realizar mediciones adicionales con una segunda bobina, $L_2$, con el objetivo de obtener resultados más coherentes para el análisis y la comparación de los distintos núcleos.

\begin{figure}[H]
\centering

\begin{circuitikz}

% Fuente
\draw
(0,0) to[sV,l=$V_S$] (0,4);

% Amperímetro
\draw (2,4) circle (0.5);
\node at (2,4) {A};

\draw (0,4) -- (1.5,4);
\draw (2.5,4) -- (3,4);

% Vatímetro
\draw (4,4) circle (0.5);
\node at (4,4) {W};

\draw (3,4) -- (3.5,4);
\draw (4.5,4) -- (10,4);

% Conexión vatímetro inferior
\draw (4,3.5) -- (4,0);

% Conexión vatímetro superior
\draw (4,4.5) -- (5,4.5) -- (5,4);

% Línea inferior
\draw (0,0) -- (10,0);

% Voltímetro
\draw (5.5,2) circle (0.5);
\node at (5.5,2) {V};

\draw (5.5,4) -- (5.5,2.5);
\draw (5.5,1.5) -- (5.5,0);

% Rama resistiva
\draw
(6.5,4) to[closing switch,l=$S_R$] (6.5,3.2)
to[R,l=$R$] (6.5,0);

% Rama inductiva
\draw
(8,4) to[closing switch,l=$S_L$] (8,3.2)
to[R,l=$R_L$] (8,2)
to[L,l=$L$] (8,0);

% Rama capacitiva
\draw
(10,4) to[closing switch,l=$S_C$] (10,3.2)
to[C,l=$C$] (10,0);

\end{circuitikz}

\caption{Circuito de análisis.}
\label{fig:circuito}

\end{figure}


\subsection{Marco teórico}

En corriente alterna, la potencia depende de los valores eficaces de tensión y corriente, y también del desfase existente entre ambas magnitudes. En este informe, $\underline{V}$ representa el fasor de tensión eficaz aplicada sobre la carga, mientras que $\underline{I}$ representa el fasor de corriente eficaz total consumida por el circuito.

La potencia aparente $S$, medida en voltampere [VA], se define como:

\begin{equation}
S = |\underline{V}|\,|\underline{I}|
\label{eq:potencia_aparente}
\end{equation}

mientras que la potencia activa $P$, medida en watt [W], está dada por:

\begin{equation}
P = |\underline{V}|\,|\underline{I}|\cos(\varphi)
\label{eq:potencia_activa}
\end{equation}

donde $\varphi$ representa el ángulo de desfase entre la tensión y la corriente. La potencia reactiva $Q$, medida en voltampere reactivo [VAr], se expresa como:

\begin{equation}
Q = |\underline{V}|\,|\underline{I}|\sin(\varphi)
\label{eq:potencia_reactiva}
\end{equation}

Estas magnitudes se relacionan mediante el triángulo de potencias:

\begin{equation}
S^2 = P^2 + Q^2
\label{eq:triangulo_potencias}
\end{equation}

El factor de potencia $fp$ es una magnitud adimensional que indica la relación entre la potencia activa y la potencia aparente del circuito. Se define como:

\begin{equation}
fp = \cos(\varphi) = \frac{P}{S}
\label{eq:factor_potencia}
\end{equation}


La impedancia equivalente $Z_{\text{eq}}$ representa la oposición total que presenta una carga o conjunto de cargas al paso de la corriente alterna. A partir de los valores eficaces de tensión y corriente, su módulo puede calcularse como:

\begin{equation}
|Z_{\text{eq}}| = \frac{|\underline{V}|}{|\underline{I}|}
\label{eq:impedancia_equivalente}
\end{equation}

donde $|Z_{\text{eq}}|$ corresponde al módulo de la impedancia equivalente vista por la fuente, $|\underline{V}|$ es el valor eficaz de la tensión aplicada sobre la carga equivalente y $|\underline{I}|$ es el valor eficaz de la corriente total consumida por dicha carga.

Además de trabajar con la impedancia equivalente, en algunos análisis resulta conveniente utilizar la admitancia equivalente del circuito. La admitancia $Y$ se define como la inversa de la impedancia $Z$:

\begin{equation}
Z=\frac{1}{Y}
\label{eq:impedancia_admitancia}
\end{equation}

A partir de la potencia compleja $\underline{S}$ y del valor eficaz de la tensión aplicada sobre la carga, la admitancia total del circuito puede calcularse como:

\begin{equation}
Y_{\text{tot}}=\frac{\underline{S}^{*}}{|\underline{V}|^2}
\label{eq:Ytot}
\end{equation}

donde $Y_{\text{tot}}$ representa la admitancia total equivalente del circuito, $\underline{S}^{*}$ es el conjugado de la potencia compleja y $|\underline{V}|$ es el valor eficaz de la tensión aplicada sobre la carga.
\newpage